% Revised Section 7 and Appendix H. Requires amsmath, amssymb, amsthm.
\section{Testing Equality and Directional Differences in VPCs}
\label{sec:equality_vpc}

The pointwise null in Section~\ref{sec:pointwise_inference} asks whether
random effects contribute variation at a specified covariate value. A
different question is whether the variance partition coefficient (VPC)
is unchanged between two values. The equality null generally permits
nonzero random effects and has a different boundary geometry. We develop likelihood-ratio tests for equality and directional
contrasts, characterize their fixed-stratum Gaussian limits, and provide
conservative compact-uniform calibrations on verified GLMM classes.
The latter allow arbitrary positive-definite Fisher-information coupling
and do not require an unproved least-favorable endpoint principle.

\subsection{Hypothesis and constrained likelihood}
\label{subsec:eq_hypothesis}
Let $x_0\ne x_1$ be fixed and write
$\theta=(\xi^{\mathsf T},\gamma^{\mathsf T})^{\mathsf T}$,
where $\xi$ includes fixed effects and all family-specific parameters
and $\gamma=\operatorname{vech}(\Gamma)$, with
$\Gamma\in\mathbb S_+^2$. Define
\begin{equation}\label{eq:eq_D}
 D(\theta)=\rho(x_1;\theta)-\rho(x_0;\theta).
\end{equation}
We test
\begin{equation}\label{eq:eq_H0}
 H_{0,E}:D(\theta)=0
 \quad\text{against}\quad H_{1,E}:D(\theta)\ne0.
\end{equation}
Let $\widehat\theta$ maximize the marginal log-likelihood over
$\Theta=\Xi\times K$, where
$K=\operatorname{vech}(\mathbb S_+^2)$, and let
$\widetilde\theta_E$ maximize it over
$\Theta_E=\{\theta\in\Theta:D(\theta)=0\}$. The statistic is
\begin{equation}\label{eq:eq_LRT}
 \mathrm{LRT}_{E,n}=2\{\ell_n(\widehat\theta)
                  -\ell_n(\widetilde\theta_E)\}.
\end{equation}
The null fit imposes the exact nonlinear equality in
\eqref{eq:eq_H0}, not merely its linearization.

For the random-intercept-and-slope design, put
$z(x)=(1,x)^{\mathsf T}$ and
$v(x;\Gamma)=z(x)^{\mathsf T}\Gamma z(x)$.
Write $\rho(x;\theta)=R_x(v(x;\Gamma),\xi)$,
where $R_x(0,\xi)=0$ and $R_x$ is differentiable at zero.
For $x_0=0$ and $x_1=1$,
\begin{equation}\label{eq:eq_D01}
\begin{aligned}
 D(\theta)={}&R_1(\gamma_{11}+2\gamma_{12}
                     +\gamma_{22},\xi)\\
 &-R_0(\gamma_{11},\xi).
\end{aligned}
\end{equation}

\subsection{Vertex null: a covariance wedge}
\label{subsec:eq_vertex_geometry}
At $\Gamma=0$, assume
\begin{equation}\label{eq:eq_rkappa}
 r_j=\partial_vR_j(0,\xi_0)>0,
 \qquad \kappa=r_0/r_1>0.
\end{equation}
Since $R_x(0,\xi)\equiv0$, the derivative of $D$ with respect
to $\xi$ vanishes at the vertex. The linearized equality is
\begin{equation}\label{eq:eq_linear_vertex}
 2H_{12}+H_{22}=(\kappa-1)H_{11}.
\end{equation}
Under the smoothness and interior-nuisance conditions
of Theorem~\ref{thm:eq_vertex_tangent_regularity},
the nonlinear equality-constrained parameter space
has the exact Bouligand tangent cone
\[
T_{\Theta_E}(\theta_0)
=
\mathbb R^{\dim\xi}
\times\operatorname{vech}(W_\kappa),
\]
where the covariance component is
\begin{equation}\label{eq:eq_wedge}
 W_\kappa=\{H\succeq0:\;
 2H_{12}+H_{22}=(\kappa-1)H_{11}\}.
\end{equation}
The tangent-cone identity is established constructively
in Appendix~\ref{app:eq_wedge_proof}, rather than
being assumed from the linearized equality alone.
With
\begin{equation}\label{eq:eq_wpm}
 w_\pm=(1,-1\pm\sqrt\kappa)^{\mathsf T},
\end{equation}
the linearized cone has the exact representation
\begin{equation}\label{eq:eq_wedge_generators}
 W_\kappa=\operatorname{cone}
 \{w_+w_+^{\mathsf T},w_-w_-^{\mathsf T}\}.
\end{equation}
Thus the vertex null has a two-dimensional covariance wedge,
rather than the single ray arising in pointwise zero-VPC testing.

\subsection{Efficient-information Gaussian limit}
\label{subsec:eq_gaussian}
Assume joint local asymptotic normality (LAN), nonsingular
normalized full Fisher information $I_0$, and the constrained
argmax and tangent-set regularity conditions in
Appendix~\ref{app:equality_proofs}. Partition $I_0$
conformably with $\theta=(\xi,\gamma)$ and define
\begin{equation}\label{eq:eq_J}
 J_0=I_{\gamma\gamma,0}
 -I_{\gamma\xi,0}I_{\xi\xi,0}^{-1}
  I_{\xi\gamma,0}.
\end{equation}
At the vertex let $B=J_0^{1/2}$, and write
$C=BK$ and $W=B\operatorname{vech}(W_\kappa)$.
For $G\sim N_3(0,I_3)$, the fixed-vertex limit is
\begin{equation}\label{eq:eq_vertex_limit}
 \mathrm{LRT}_{E,n}\ \Longrightarrow\
 T_0:=\|\Pi_C(G)\|^2-\|\Pi_W(G)\|^2.
\end{equation}
For a local null sequence with
$\sqrt n\,B\gamma_n\to\mu\in W$ and a common limiting
information matrix, the corresponding Gaussian limit is
\begin{equation}\label{eq:eq_local_limit}
 T_\mu=\|\Pi_C(G+\mu)\|^2
       -\|\Pi_W(G+\mu)\|^2.
\end{equation}
This shifted expression describes vertex-local null sequences,
subject to the stated moving-parameter regularity conditions;
it does not describe every fixed nonzero covariance matrix.

\subsection{Exact spherical-angular law at the vertex}
\label{subsec:eq_vertex_angular}
The fixed-vertex distribution in \eqref{eq:eq_vertex_limit}
can be evaluated by deterministic angular integration. Define,
for $u\in S^2=\{u\in\mathbb R^3:\|u\|=1\}$,
\begin{equation}\label{eq:eq_vertex_angular_function}
 d_E(u)=\|\Pi_C(u)\|^2-\|\Pi_W(u)\|^2\ge0.
\end{equation}
The inequality follows from $W\subseteq C$.
Let $\overline F_{\chi_3^2}$ denote the survival function
of a chi-square variable with three degrees of freedom.
Then, for $t>0$,
\begin{equation}\label{eq:eq_vertex_angular_tail}
 \boxed{\displaystyle
 \Pr(T_0>t)=\frac{1}{4\pi}
 \int_{\{u\in S^2:d_E(u)>0\}}
 \overline F_{\chi_3^2}\!\left(\frac{t}{d_E(u)}\right)
 \,d\omega(u).}
\end{equation}
Equivalently, writing
$u(\theta,\phi)=(\sin\phi\cos\theta,
\sin\phi\sin\theta,\cos\phi)^{\mathsf T}$,
\begin{equation}\label{eq:eq_vertex_spherical_tail}
\begin{aligned}
 \Pr(T_0>t)=\frac{1}{4\pi}
 \int_0^{2\pi}\!\int_0^\pi
 &\mathbf1_{\{d_E(u(\theta,\phi))>0\}}\\
 &\times\overline F_{\chi_3^2}\!\left(
 \frac{t}{d_E(u(\theta,\phi))}\right)
 \sin\phi\,d\phi\,d\theta.
\end{aligned}
\end{equation}
The distribution may have an atom at zero, with mass
\begin{equation}\label{eq:eq_vertex_atom}
 \Pr(T_0=0)=\frac{1}{4\pi}
 \int_{S^2}\mathbf1_{\{d_E(u)=0\}}\,d\omega(u).
\end{equation}
For $t>0$ its density is
\begin{equation}\label{eq:eq_vertex_density}
 f_{T_0}(t)=\frac{1}{4\pi}
 \int_{\{d_E(u)>0\}}
 \frac{1}{d_E(u)}
 f_{\chi_3^2}\!\left(\frac{t}{d_E(u)}\right)
 \,d\omega(u).
\end{equation}
Appendix~\ref{app:eq_vertex_angular_proof} proves these
identities and gives an explicit two-generator projection
algorithm for $W$. Thus the fixed-vertex tail and critical
values can be computed by spherical quadrature, without
simulation of Gaussian vectors. This exact fixed-vertex law
must not be substituted for the uniform Gaussian-local bound
in Section~\ref{subsec:eq_calibration}: the latter covers
all null displacements $\mu\in W$, not only $\mu=0$.

\subsection{Nonzero rank-one null and angular law}
\label{subsec:eq_rankone}
Suppose $\Gamma_0$ has rank one and a positive eigenvalue,
$D(\theta_0)=0$, and the full-parameter gradient
$a=\nabla_\theta D(\theta_0)$ is nonzero. Let $n$ denote the
full-parameter inward normal to the active PSD boundary,
so the alternative tangent is $\{h:n^{\mathsf T}h\ge0\}$.
Assume the null tangent is exactly
\begin{equation}\label{eq:eq_rankone_tangent}
 \{h:n^{\mathsf T}h\ge0,\ a^{\mathsf T}h=0\}.
\end{equation}
Define the normalized Fisher-whitened normal inner product
\begin{equation}\label{eq:eq_angle}
 c=\frac{a^{\mathsf T}I_0^{-1}n}
 {\sqrt{(a^{\mathsf T}I_0^{-1}a)
             (n^{\mathsf T}I_0^{-1}n)}}.
\end{equation}
When $|c|<1$, whitening and removal of the common lineality
space reduce the nested tangent-cone problem to a half-plane
versus a ray in two dimensions. Choose coordinates with
$H=\{(z_1,z_2):z_2\ge0\}$ and the whitened equality normal
$(\sqrt{1-c^2},c)^{\mathsf T}$. The equality line intersects
$H$ in the ray generated by
$(-c,\sqrt{1-c^2})^{\mathsf T}$.
Consequently its angle from the positive horizontal axis is
$\omega=\arccos(-c)\in(0,\pi)$.
For independent $Z_1,Z_2\sim N(0,1)$, the limiting statistic is
\begin{equation}\label{eq:eq_angular_law}
\begin{aligned}
 T_\omega={}&Z_1^2+(Z_2)_+^2\\
 &-(Z_1\cos\omega+Z_2\sin\omega)_+^2,
\end{aligned}
\end{equation}
where $y_+=\max(y,0)$. Reflection symmetry implies that
$\omega$ and $\pi-\omega$ give the same law. When $c=0$,
$T_\omega\sim\chi_1^2$.

For $t>0$, define
\begin{equation}\label{eq:eq_homega}
 h_\omega(\phi)=\cos^2\phi+(\sin\phi)_+^2
                  -(\cos(\phi-\omega))_+^2.
\end{equation}
The upper-tail probability is
\begin{equation}\label{eq:eq_angular_tail}
 \Pr(T_\omega>t)=\frac1{2\pi}
 \int_{\{\phi\in[0,2\pi):h_\omega(\phi)>0\}}
 \exp\!\left\{-\frac{t}{2h_\omega(\phi)}\right\}
 \,d\phi.
\end{equation}
At $|c|=1$, the first-order tangent description may be
singular; the boundary-normal formula alone does not justify
an LRT limit without further analysis.

\subsection{Conservative Gaussian calibration under arbitrary information coupling}
\label{subsec:eq_calibration}
We now construct coordinates that simultaneously preserve the full
PSD alternative and put the linearized equality null in canonical form.
Let $P_\kappa=[w_+\ \ w_-]$, with $w_\pm$ as in
\eqref{eq:eq_wpm}. Since $\det P_\kappa=-2\sqrt\kappa\ne0$,
the congruence map
\begin{equation}\label{eq:eq_canonical_map}
 \Gamma=P_\kappa
 \begin{pmatrix}a&z\\z&b\end{pmatrix}P_\kappa^{\mathsf T}
\end{equation}
is an invertible linear reparameterization of the covariance
coordinates. Congruence preserves positive semidefiniteness, and
substitution into the vertex contrast derivative
\eqref{eq:eq_linear_vertex} shows that the linearized equality is
exactly $z=0$. Thus the transformed alternative and null are
\begin{equation}\label{eq:eq_canonical_cones}
 \begin{aligned}
 C_0&=\{(a,b,z):a\ge0,\ b\ge0,\ ab\ge z^2\},\\
 W_0&=\{(a,b,0):a,b\ge0\}.
 \end{aligned}
\end{equation}
Define the invertible $3\times3$ matrix $L_\kappa$ by
$\operatorname{vech}(\Gamma)=L_\kappa(a,b,z)^{\mathsf T}$,
using the same ordering of $\operatorname{vech}$ as in $J_0$.
The information in canonical coordinates is
\begin{equation}\label{eq:eq_canonical_Q}
 Q=L_\kappa^{\mathsf T}J_0L_\kappa\succ0.
\end{equation}
This is a congruence transformation of the efficient information,
not the untransformed matrix $J_0$. In particular, the bounds below
apply to the precise cone pair in \eqref{eq:eq_canonical_cones},
not to an arbitrary image of the PSD cone.
For $X\sim N_3((a,b,0),Q^{-1})$ define
\begin{equation}\label{eq:eq_Qstat}
 \begin{aligned}
 T_Q(X)&=d_Q^2(X,W_0)-d_Q^2(X,C_0),\\
 F_{Q,t}(a,b)&=\Pr\{T_Q(X)>t\}.
 \end{aligned}
\end{equation}
The general Gaussian-local worst-case survival probability is
$S_Q^\star(t)=\sup_{a,b\geq0}F_{Q,t}(a,b)$.

\paragraph{A baseline distance bound.}
For $W\subseteq C$ closed convex cones and $\mu\in W$,
\begin{equation}\label{eq:eq_pathwise_bound}
 0\leq T_\mu\leq d^2(G,W)\quad\text{almost surely},\qquad G\sim N_3(0,I_3).
\end{equation}
For a planar wedge with opening angle $\vartheta$ this distance has
survival function
\begin{equation}\label{eq:eq_old_bound}
 B_Q^{\rm old}(t)=\frac{\vartheta}{2\pi}\overline F_{\chi_1^2}(t)
 +\frac12\overline F_{\chi_2^2}(t)
 +\frac{\pi-\vartheta}{2\pi}\overline F_{\chi_3^2}(t).
\end{equation}
Here $\vartheta=\arccos\{Q_{ab}/\sqrt{Q_{aa}Q_{bb}}\}$,
which is the angle between the two null generators in the $Q$ metric.

\paragraph{A sharper coupling-sensitive envelope.}
Partition the transformed information as
\begin{equation}\label{eq:eq_Qblocks}
 \begin{aligned}
 Q&=\begin{pmatrix}A&h\\h^{\mathsf T}&q\end{pmatrix}\succ0,\\
 D&=q-h^{\mathsf T}A^{-1}h>0,\\
 \kappa_Q^2&=\frac{h^{\mathsf T}A^{-1}h}{D},
 \quad L=\kappa_Q.
 \end{aligned}
\end{equation}
The subscript on $\kappa_Q$ distinguishes information coupling
from the VPC derivative ratio $\kappa=r_0/r_1$. The latter is a
local VPC derivative ratio, not the CPG family power parameter.
Set $K_A=A^{1/2}\mathbb R_+^2$ and
$d=A^{-1/2}h/\sqrt D$; define the orientation constants
$\lambda_+=d(d,K_A)$ and $\lambda_-=d(-d,K_A)$, where $d(\cdot,K_A)$
denotes Euclidean distance. Let $V\sim N(0,1)$ and, independently,
$R^2\sim w_0\chi_0^2+w_1\chi_1^2+w_2\chi_2^2$ with
\begin{equation}\label{eq:eq_wmix}
 (w_0,w_1,w_2)=\left(\frac{\vartheta}{2\pi},\frac12,
 \frac{\pi-\vartheta}{2\pi}\right).
\end{equation}
For $v,r\geq0$ and $0\leq\lambda\leq L$, put
\begin{equation}\label{eq:eq_Blambda}
 \begin{aligned}
 B_{L,\lambda}(v,r)&=v^2+r^2-\frac{(r-\lambda v)_+^2}{1+L^2},\\
 J_Q(t)&=\frac12\sum_{\epsilon\in\{+,-\}}
 \Pr\{B_{L,\lambda_\epsilon}(|V|,R)>t\}.
 \end{aligned}
\end{equation}
Define $p_1(t)=\overline F_{\chi_1^2}(t)$ and $M=1+\kappa_Q^2$.
A complementary closed-form envelope is
\begin{equation}\label{eq:eq_Hbound}
 \begin{split}
 H_Q(t)={}&\frac12p_1(t)+\frac12p_1(t/M)\\
 &+\frac{\pi-\vartheta}{2\pi}\frac{\kappa_Q}{\sqrt M}
 \exp\!\left(-\frac{t}{2M}\right)
 \left\{2\Phi\!\left(\frac{\kappa_Q\sqrt t}{\sqrt M}\right)-1\right\}.
 \end{split}
\end{equation}
\begin{theorem}[Improved Gaussian-local calibration]\label{thm:eq_improved}
For every fixed $Q\succ0$, $t>0$, and $(a,b)\in\mathbb R_+^2$,
\begin{equation}\label{eq:eq_improved}
 F_{Q,t}(a,b)\leq \mathcal U_Q(t):=\min\{J_Q(t),H_Q(t)\}
 <B_Q^{\rm old}(t).
\end{equation}
The function $\mathcal U_Q$ is continuous and strictly decreasing from
one to zero on $t\geq0$. Thus for $\alpha\in(0,1)$ the unique positive
solution $c_\alpha(Q)$ of $\mathcal U_Q(c_\alpha(Q))=\alpha$
satisfies $\sup_{a,b\geq0}\Pr\{T_Q>c_\alpha(Q)\}\leq\alpha$.
The critical value is continuous in $Q\succ0$.
\end{theorem}
For this \emph{vertex-local linear chart}, $Q$ is constructed from
the vertex information and derivative ratio $\\kappa$. This is not a
valid prescription for evaluating the same information matrix at an
arbitrary nonzero covariance. The exact nuisance-dependent chart
introduced below gives the appropriate transformed metric across ranks.
Theorem~\ref{thm:eq_improved} is a result for the Gaussian experiment.
Its proof, including the orientation refinement and the closed-form bound,
is in Appendix~\ref{app:eq_improved_proof}. The general four-endpoint
identity is not assumed. The two upper bounds $J_Q$ and $H_Q$ are
separately valid, but only the orientation bound $J_Q$ is guaranteed
to improve $B_Q^{\rm old}$ strictly; $H_Q$ need not do so on its own.
The minimum in \eqref{eq:eq_improved} retains the stated improvement.

\paragraph{Exact information subclasses and the endpoint question.}
Define $E_Q(t)$ as the maximum of the vertex, two rank-one escape,
and interior escape survival limits. The rank-one endpoint theorem and
Schur-complement limits identify these four quantities. The equality
$S_Q^\star(t)=E_Q(t)$ is established if $h=0$ (the covariance contrast
is information-orthogonal to both variance coordinates) and also if
$Q_{ab}=Q_{bz}=0$, or the swapped case $Q_{ab}=Q_{az}=0$.
In the latter class, writing $s=-Q_{az}/\sqrt{Q_{aa}Q_{zz}}$,
\begin{equation}\label{eq:eq_isolated_exact}
 S_Q^\star(t)=\max\{S_R(t;0,s),\overline F_{\chi_1^2}(t)\}.
\end{equation}
The corresponding exact-case arguments are summarized in
Appendix~\ref{app:eq_exact_classes}.
For arbitrary coupled $Q$, however, the finite-displacement inequality
$F_{Q,t}(a,b)\leq E_Q(t)$ remains an \emph{open conjecture}.
The escape theorem excludes an obstruction at infinity but not at a
finite interior or boundary-ray displacement.

\paragraph{From the Gaussian experiment to the GLMM.}
At a fixed vertex-local null sequence, convergence of the joint
likelihood-ratio statistic and a consistent $Q$ estimator transfers
this Gaussian bound through continuity at its positive critical value.
That statement alone is not uniform across covariance ranks. The
exact nuisance-dependent chart and the uniform likelihood argument in
Subsection~\ref{subsec:sec7-uniform-contrast} supply the corresponding
compact-uniform conclusion on the verified four-family classes.

\subsection{Implementation and scope of validity}
\label{subsec:eq_implementation}
Implementation fits the unrestricted and exact nonlinearly
constrained marginal likelihoods and computes
\eqref{eq:eq_LRT}. At a fixed vertex, the efficient information
and wedge generators determine the exact angular law
\eqref{eq:eq_vertex_angular_tail}; deterministic spherical
quadrature provides its tail probabilities and quantiles.
For conservative calibration over vertex-local null shifts,
compute the transformed information $Q$ and use the sharper
critical value from Theorem~\ref{thm:eq_improved}.
At a fixed nonzero rank-one null, the angular law
\eqref{eq:eq_angular_law} is the appropriate fixed-parameter
limit under the stated transversality and tangent regularity
conditions. At a fixed positive-definite null, the regular
equality LRT has a $\chi_1^2$ limit when $\nabla D\ne0$.

Theorem~\ref{thm:eq_improved} controls the Gaussian experiment.
A single compact-uniform GLMM rule is given below in
Theorem~\ref{thm:sec7-uniform-contrasts}; it uses the
\emph{full-information-transformed} metric at the true covariance,
not the original vertex-information plug-in at every rank.
The arbitrary-coupling four-endpoint identity remains open as a
question of sharper calibration, not of conservative test validity.


\subsection{Directional inference for VPC differences}
\label{sec:directional-vpc-difference}

The equality test developed in Sections 7.1--7.6 assesses
whether the VPC changes between two covariate values.
In applications, the direction of that change may itself be
of scientific interest. We therefore consider the one-sided
hypotheses
\begin{equation}
H_{0,+}:D(\theta)\leq 0
\qquad\text{versus}\qquad
H_{1,+}:D(\theta)>0,
\label{eq:directional-hypotheses}
\end{equation}
where
\begin{equation*}
D(\theta)
=
\rho(x_1;\theta)-\rho(x_0;\theta).
\end{equation*}
The reverse-direction test interchanges the inequalities.

Let
\begin{equation*}
\Theta_-
=
\left\{
\theta\in\Theta:
D(\theta)\leq 0
\right\},
\end{equation*}
and let $\widehat{\theta}_-$ maximize the marginal
log-likelihood over $\Theta_-$. With $\widehat{\theta}$
denoting the unrestricted constrained MLE over
$\Theta=\Xi\times K$, define
\begin{equation}
\mathrm{LRT}_{+,n}
=
2\left\{
\ell_n(\widehat{\theta})
-
\ell_n(\widehat{\theta}_-)
\right\}.
\label{eq:directional-lrt}
\end{equation}
The restricted likelihood imposes the exact nonlinear
inequality $D(\theta)\leq 0$. The linearized constraints
below are used only to characterize local asymptotic
distributions.

\paragraph{Vertex null geometry.}

Consider $x_0=0$ and $x_1=1$, and let
$\theta_*=(\xi_*^\top,0^\top)^\top$ be a covariance-vertex
parameter. Retain the notation
\begin{equation*}
r_j
=
\partial_v R_j(0,\xi_*)>0,
\qquad
\kappa
=
\frac{r_0}{r_1}
\end{equation*}
from Section 7.2. By Appendix~\ref{app:eq_wedge_proof}, the first-order
contrast derivative along a feasible covariance
direction $H$ is
\begin{equation}
\dot{D}_{\theta_*}(0,H)
=
r_1
\left\{
2H_{12}+H_{22}
-
(\kappa-1)H_{11}
\right\}.
\label{eq:directional-vertex-derivative}
\end{equation}

Define the directional covariance cones
\begin{align}
K_-(\kappa)
&=
\left\{
H\succeq 0:
2H_{12}+H_{22}
\leq
(\kappa-1)H_{11}
\right\},
\label{eq:directional-negative-cone}
\\
K_+(\kappa)
&=
\left\{
H\succeq 0:
2H_{12}+H_{22}
\geq
(\kappa-1)H_{11}
\right\}.
\label{eq:directional-positive-cone}
\end{align}
These are closed convex cones satisfying
\begin{align}
K_-(\kappa)\cup K_+(\kappa)
&=
\mathbb{S}_+^2,
\\
K_-(\kappa)\cap K_+(\kappa)
&=
W_\kappa,
\end{align}
where $W_\kappa$ is the equality wedge in \eqref{eq:eq_wedge}.

The identification of $K_-(\kappa)$ as the covariance
tangent cone to the exact directional null requires
tangential regularity. The next proposition states the generic
fixed-parameter limit under that condition. For the four
strictly monotone VPC links, the exact chart in
Subsection~\ref{subsec:sec7-exact-chart} verifies this geometry,
including at the covariance vertex.

\begin{proposition}[Directional vertex limit]
\label{prop:directional-vertex-limit}

Suppose the joint LAN, information nonsingularity,
constrained-argmax, and local-set convergence conditions
of Section 4 hold at $\theta_*$. Assume additionally
that the covariance tangent cone of $\Theta_-$ at
$\theta_*$ is $K_-(\kappa)$, after removal of the
unrestricted nuisance coordinates.

Let $J_*$ be the efficient covariance information and
define
\begin{equation*}
C
=
J_*^{1/2}K,
\qquad
C_-
=
J_*^{1/2}
\operatorname{vech}
\left\{
K_-(\kappa)
\right\}.
\end{equation*}

Then
\begin{equation}
\mathrm{LRT}_{+,n}
\xrightarrow{d}
T_{+,0}
=
\left\|
\Pi_C(G)
\right\|^2
-
\left\|
\Pi_{C_-}(G)
\right\|^2,
\qquad
G\sim N_3(0,I_3).
\label{eq:directional-vertex-limit}
\end{equation}

For local directional-null sequences satisfying
\begin{equation*}
\sqrt{n}\,
J_*^{1/2}\gamma_n
\longrightarrow
\mu\in C_-,
\end{equation*}
under the corresponding joint Gaussian-shift and
local-set conditions, the limiting statistic is
\begin{equation}
T_{+,\mu}
=
\left\|
\Pi_C(G+\mu)
\right\|^2
-
\left\|
\Pi_{C_-}(G+\mu)
\right\|^2.
\label{eq:directional-shift-limit}
\end{equation}

\end{proposition}

\paragraph{Conservative Gaussian calibration.}

The distance-bound theorem of Appendix H applies directly to the nested
closed convex cones $C_-\subseteq C$. Consequently,
for every $\mu\in C_-$,
\begin{equation}
0
\leq
T_{+,\mu}
\leq
d^2(G,C_-)
\qquad\text{almost surely}.
\label{eq:directional-distance-bound}
\end{equation}
In particular,
\begin{equation}
\sup_{\mu\in C_-}
\Pr(T_{+,\mu}>t)
\leq
\Pr
\left\{
d^2(G,C_-)>t
\right\}.
\label{eq:directional-tail-bound}
\end{equation}

By Moreau decomposition,
\begin{equation*}
d^2(G,C_-)
=
\left\|
\Pi_{C_-^\circ}(G)
\right\|^2.
\end{equation*}
Thus its distribution is the chi-bar-square mixture
\begin{equation}
d^2(G,C_-)
\sim
\sum_{j=0}^{3}
v_j(C_-^\circ)\chi_j^2,
\label{eq:directional-distance-mixture}
\end{equation}
where the sum denotes a distributional mixture and
$v_j(C_-^\circ)$ are the conic intrinsic volumes of the
polar cone.

Let $c_{\alpha,-}$ denote an upper-tail critical value
satisfying
\begin{equation*}
\Pr
\left\{
d^2(G,C_-)>c_{\alpha,-}
\right\}
\leq
\alpha.
\end{equation*}
Then
\begin{equation}
\sup_{\mu\in C_-}
\Pr
\left(
T_{+,\mu}>c_{\alpha,-}
\right)
\leq
\alpha.
\label{eq:directional-gaussian-calibration}
\end{equation}

This bound is uniform over the indicated Gaussian
local-null shifts. It does not, without additional
assumptions, establish uniform size control for the
original GLMM over all covariance ranks.

\paragraph{Nonzero rank-one equality points.}

Suppose $\Gamma_0$ has rank one with a positive
eigenvalue, $D(\theta_0)=0$, and the full-parameter
contrast gradient
\begin{equation*}
a
=
\nabla_\theta D(\theta_0)
\end{equation*}
is nonzero.

Let $n$ be the supporting normal to the active PSD
boundary, oriented so that the alternative tangent is
\begin{equation*}
T_1
=
\left\{
h:
n^\top h\geq 0
\right\}.
\end{equation*}
Assume that the exact directional-null tangent is
\begin{equation*}
T_-
=
\left\{
h:
n^\top h\geq 0,
\quad
a^\top h\leq 0
\right\}.
\end{equation*}

Retain the signed information-metric correlation
\begin{equation}
c
=
\frac{
a^\top I_0^{-1}n
}{
\sqrt{
(a^\top I_0^{-1}a)
(n^\top I_0^{-1}n)
}
},
\label{eq:directional-normal-correlation}
\end{equation}
and suppose $|c|<1$. In the coordinates below, the
inward PSD normal points along the positive $v$ axis, so the
unit contrast normal has components $(\sqrt{1-c^2},c)$.

\begin{proposition}[Directional rank-one limit]
\label{prop:directional-rankone-limit}

Under the preceding conditions and the joint
constrained likelihood regularity assumptions of
Section 4, the directional LRT has the limiting
distribution
\begin{equation}
\mathrm{LRT}_{+,n}
\xrightarrow{d}
T_{+,c}
=
\left\|
\Pi_{\mathcal H}(Z)
\right\|^2
-
\left\|
\Pi_{\mathcal Q_c}(Z)
\right\|^2,
\label{eq:directional-rankone-limit}
\end{equation}
where
\begin{equation*}
Z\sim N_2(0,I_2),
\end{equation*}
and
\begin{align*}
\mathcal H
&=
\left\{
(u,v):
v\geq 0
\right\},
\\
\mathcal Q_c
&=
\left\{
(u,v):
v\geq 0,
\quad
\sqrt{1-c^2}\,u+cv\leq 0
\right\}.
\end{align*}

For
\begin{equation*}
u_\varphi
=
(\cos\varphi,\sin\varphi)^\top,
\end{equation*}
define
\begin{equation*}
h_c(\varphi)
=
\left\|
\Pi_{\mathcal H}(u_\varphi)
\right\|^2
-
\left\|
\Pi_{\mathcal Q_c}(u_\varphi)
\right\|^2.
\end{equation*}

Then, for $t>0$,
\begin{equation}
\Pr(T_{+,c}>t)
=
\frac{1}{2\pi}
\int_{\{h_c(\varphi)>0\}}
\exp
\left\{
-\frac{t}{2h_c(\varphi)}
\right\}
\,d\varphi.
\label{eq:directional-angular-tail}
\end{equation}

\end{proposition}

In a general model, the cases $|c|=1$ require separate analysis.
For the four strictly increasing variance links here, the equality
covariance gradient is an indefinite matrix whereas the rank-one
PSD boundary normal is semidefinite. Thus these two normals are
linearly independent and $|c|<1$ at every regular nonzero rank-one
equality point. This is established in the completed geometric proof
and applies to the directional boundary as well.

\paragraph{Regular and strictly null configurations.}

At a fixed positive-definite equality point with
$\nabla_\theta D(\theta_0)\neq 0$, standard one-sided
likelihood theory yields
\begin{equation}
\mathrm{LRT}_{+,n}
\xrightarrow{d}
\frac{1}{2}\chi_0^2
+
\frac{1}{2}\chi_1^2.
\label{eq:directional-interior-limit}
\end{equation}

At a fixed parameter satisfying $D(\theta_0)<0$,
the directional inequality is inactive locally.
Under consistency of the unrestricted and restricted
estimators and the usual likelihood regularity
conditions,
\begin{equation}
\mathrm{LRT}_{+,n}
\xrightarrow{p}
0.
\label{eq:directional-strict-null-limit}
\end{equation}

The vertex Gaussian calibration, fixed rank-one angular law,
and regular interior limit characterize distinct regimes.
The following exact-coordinate construction and uniform transfer
result combine them into a conservatively calibrated test over
the specified compact informative-design classes, including
covariance-rank transitions.


\subsection{Exact PSD-preserving contrast coordinates}
\label{subsec:sec7-exact-chart}

The preceding linear canonical chart describes the vertex geometry,
but the VPC equality constraint is generally nonlinear away from the
vertex. An exact local change of covariance coordinates is possible
because the four variance-link functions are strictly increasing.
Let $z_j=(1,x_j)^{\mathsf T}$, $j=0,1$, with $x_0\ne x_1$,
put $B=[z_0\ z_1]$, and write
\begin{equation}\label{eq:sec7-twolevel-cov}
 B^{\mathsf T}\Gamma B=\begin{pmatrix}v_0&c\\c&v_1\end{pmatrix}=:M.
\end{equation}
The map $\Gamma\mapsto M$ is invertible and preserves PSD.
Retain $R_j(v,\xi)=\rho(x_j;\xi,\Gamma)$ at $v=v_j$.
On the relevant compact directional-null class and a neighborhood
thereof, define the smooth inverse-link map
\begin{equation}\label{eq:sec7-hlink}
 h_\xi(v)=R_0(\cdot,\xi)^{-1}\{R_1(v,\xi)\},\qquad
 s_\xi(v)=\begin{cases}
 h_\xi(v)/v,&v>0,\\
 h_\xi'(0),&v=0.
 \end{cases}
\end{equation}
Here $h_\xi(0)=0$ and
$h_\xi'(0)=R_{1,v}(0,\xi)/R_{0,v}(0,\xi)>0$.
In particular $s_\xi$ has a smooth positive extension through zero.

\begin{theorem}[Exact contrast chart]\label{thm:sec7-exact-chart}
Suppose $R_j(0,\xi)=0$ and $R_j$ is $C^5$ and strictly
increasing in variance, with derivative bounded above and away
from zero on the relevant buffered compact class. These properties
hold for the Gaussian identity-link and the Poisson, NB2, and CPG
log-link VPCs under their admissible compact family restrictions.
There is a nuisance-preserving, locally invertible $C^3$ change
of coordinates near every compact equality or directional null
class, given by
\begin{equation}\label{eq:sec7-exact-psd-map}
 \begin{aligned}
 H_\xi(M)&=
 \begin{pmatrix}
 v_0 &c\sqrt{s_\xi(v_1)}\\
 c\sqrt{s_\xi(v_1)}&h_\xi(v_1)
 \end{pmatrix},\\
 \widetilde\Gamma&=O^{\mathsf T}H_\xi(M)O,
 \quad O=\frac1{\sqrt2}
 \begin{pmatrix}1&1\\1&-1\end{pmatrix}.
 \end{aligned}
\end{equation}
It preserves positive semidefiniteness and covariance rank and
has a smooth local inverse. With
$\nu=(\widetilde\gamma_{11},\widetilde\gamma_{12},
\widetilde\gamma_{22})^{\mathsf T}$, the \emph{exact} hypotheses are
\begin{equation}\label{eq:sec7-chart-null-cones}
 \begin{aligned}
 D=0&\iff\widetilde\gamma_{12}=0,\\
 D\le0&\iff\widetilde\gamma_{12}\ge0.
 \end{aligned}
\end{equation}
Consequently, locally on the buffered true null classes, the
alternative, equality, and directional covariance factors are
respectively the fixed closed convex cones
\begin{equation}\label{eq:sec7-fixed-cones}
 \begin{aligned}
 K&=\{\nu:\widetilde\Gamma(\nu)\succeq0\},\\
 W_E&=K\cap\{\nu:\nu_2=0\},\\
 W_-&=K\cap\{\nu:\nu_2\ge0\}.
 \end{aligned}
\end{equation}
\end{theorem}

\begin{proof}
The strict monotonicity of $R_j$ ensures that $h_\xi$ exists
on a common neighborhood of the compact null class: on the
directional null,
$R_1(v_1,\xi)\le R_0(v_0,\xi)$, so
$0\le h_\xi(v_1)\le v_0$. The same inequality for every
$0\le v\le v_1$ yields a compact downward domain on which
the implicit-function theorem provides smooth inverse branches.
The identity
$s_\xi(v)=\int_0^1h_\xi'(tv)\,dt$ proves its positive smooth
extension at zero. Writing
$D_s=\operatorname{diag}(1,\sqrt{s_\xi(v_1)})$ gives
$H_\xi(M)=D_s M D_s$; hence PSD and rank are preserved in
both directions. The Jacobian of
$(v_0,v_1,c)\mapsto(v_0,h_\xi(v_1),c\sqrt{s_\xi(v_1)})$
has strictly positive determinant
$h_\xi'(v_1)\sqrt{s_\xi(v_1)}$ and an explicit local inverse.
Finally the off-diagonal element of $O^{\mathsf T}H_\xi(M)O$
is $\{v_0-h_\xi(v_1)\}/2$. Because $R_0$ is increasing,
$\operatorname{sign}D=\operatorname{sign}\{h_\xi(v_1)-v_0\}$,
which proves \eqref{eq:sec7-chart-null-cones} and the
fixed-cone representation.
\end{proof}

\paragraph{Correct information under the nonlinear chart.}
The chart depends on nuisance parameters. If
$\eta=(\xi^{\mathsf T},\nu^{\mathsf T})^{\mathsf T}$ and
$F(\eta)=(\xi^{\mathsf T},\gamma(\xi,\nu)^{\mathsf T})^{\mathsf T}$
is its inverse, set $L(\eta)=D_\eta F(\eta)$ and transform the
\emph{full} expected Fisher information before profiling:
\begin{equation}\label{eq:sec7-transform-full-info}
 \begin{aligned}
 \mathcal I(\eta)&=L(\eta)^{\mathsf T}I\{F(\eta)\}L(\eta),\\
 \mathcal J(\eta)&=\mathcal I_{\nu\nu}-
 \mathcal I_{\nu\xi}\mathcal I_{\xi\xi}^{-1}
 \mathcal I_{\xi\nu}.
 \end{aligned}
\end{equation}
At nonzero covariance, $D_\xi\gamma(\xi,\nu)$ need not vanish.
Therefore, congruence of the original covariance Schur complement
alone is generally incorrect. At the vertex, that derivative is zero,
and the ordinary covariance-only congruence is recovered.
For consistency with the Gaussian wedge formulas above, let
$P$ be the permutation satisfying
$\nu=P(a,b,z)^{\mathsf T}=(a,z,b)^{\mathsf T}$ and put
\begin{equation}\label{eq:sec7-chart-Q}
 Q=P^{\mathsf T}\mathcal J P\succ0.
\end{equation}
This $Q$ is the information metric required by the exact
nonlinear chart; it may have arbitrary off-diagonal coupling.

\subsection{Compact-uniform conservative comparison tests}
\label{subsec:sec7-uniform-contrast}

The following result strengthens the fixed-stratum Gaussian
statements to a single conservative uniform test on the verified
four-family parameter classes. It does not assert unrestricted
uniformity over unbounded parameters, uninformative designs, or
singular Fisher information.

\begin{theorem}[Uniform contrast calibration]
\label{thm:sec7-uniform-contrasts}
Suppose the finite-type informative designs, nuisance compactness,
buffered covariance domain, and exact global marginal-likelihood
conditions of Subsection~\ref{subsec:sec4-compact-uniform} hold.
The nuisance dimension is fixed and finite, the full Fisher
information is uniformly positive definite, and the family-specific
likelihood regularity conditions are as verified for the Gaussian,
Poisson, NB2, and CPG models; both CPG dispersion and power remain
unknown. Let $\mathcal N_E\subset\Theta_E$ and
$\mathcal N_-\subset\Theta_-$ be buffered compact true null
classes, possibly containing every covariance rank. For
$0<\alpha<1$, let $\widehat Q_n$ consistently estimate
\eqref{eq:sec7-chart-Q} using the correctly transformed full
information at a consistent unrestricted or corresponding null fit.
Then the exact equality LRT obeys
\begin{equation}\label{eq:sec7-equality-uniform}
 \boxed{\displaystyle
 \limsup_{n\to\infty}\sup_{\theta\in\mathcal N_E}
 P_\theta\!\left\{\mathrm{LRT}_{E,n}>
 c_\alpha(\widehat Q_n)\right\}\le\alpha,}
\end{equation}
where $c_\alpha$ is the positive conservative Gaussian
critical value from Theorem~\ref{thm:eq_improved}.
For the directional test, let $c_{\alpha,-}^{\rm dist}(Q)$ be
the positive upper-$\alpha$ quantile of the wedge-distance law
\eqref{eq:eq_old_bound}, with opening angle computed from the
correct $Q$. Then
\begin{equation}\label{eq:sec7-directional-uniform}
 \boxed{\displaystyle
 \limsup_{n\to\infty}\sup_{\theta\in\mathcal N_-}
 P_\theta\!\left\{\mathrm{LRT}_{+,n}>
 c_{\alpha,-}^{\rm dist}(\widehat Q_n)\right\}\le\alpha.}
\end{equation}
These statements allow arbitrary SPD information coupling,
rank transitions, and sequences of order $n^{-1/2}$ or larger
in covariance coordinates. They make no four-endpoint
least-favorable assertion.
\end{theorem}

\begin{proof}
In the exact chart, the local alternative and null covariance
sets centered at a true transformed covariance $\nu_n$ are
$K-\sqrt n\nu_n$ and $W_0-\sqrt n\nu_n$, with
$W_0=W_E$ or $W_-$. They are nested closed convex sets
containing zero, including along divergent local displacements.
The verified compact likelihood regularity gives uniform
root-$n$ localization of both constrained maxima, a
true-parameter-centered quadratic likelihood expansion, a
uniform Gaussian score approximation, and consistently estimated
transformed information. Profiling nuisance in the quadratic
experiment yields the corresponding Fisher-metric difference
of squared distances. Nonexpansiveness of convex projection
and positive-threshold anti-concentration of nested distance
statistics transfer Gaussian rejection probabilities uniformly,
without assuming a fixed covariance rank.

For equality the mean displacement belongs to $W_E$, so the
Gaussian envelope from Theorem~\ref{thm:eq_improved} bounds
its survival probability for every displacement and SPD $Q$.
For the directional null, the translated-cone argument gives
\[T_{-,\mu}\le d^2(G,\mathcal J^{1/2}W_-)
\le d^2(G,\mathcal J^{1/2}W_E),\]
where $G\sim N_3(0,I_3)$ and the last squared distance has
the wedge-distance law \eqref{eq:eq_old_bound}. Continuity of
the positive critical values and consistency of $\widehat Q_n$
finish both uniform size bounds. Full details of the moving-set
anti-concentration and family-specific likelihood arguments
are supplied in the completed proof supplement, Sections 2--4
and 9.
\end{proof}

\begin{corollary}[Information-free comparison rule]
\label{cor:sec7-information-free}
Under the same compact-uniform assumptions, both the equality
and directional LRTs are asymptotically level $\alpha$ uniformly
when calibrated by the fixed upper-$\alpha$ quantile of the
distributional mixture
\begin{equation}\label{eq:sec7-universal-mixture}
 \boxed{\frac12\chi_2^2+\frac12\chi_3^2.}
\end{equation}
No estimation of the chart information is required for the
critical value, although positive-definite information and the
exact global likelihood assumptions remain necessary.
\end{corollary}
\begin{proof}
For any SPD $Q$, the wedge-distance survival in
\eqref{eq:eq_old_bound} is bounded by
$\tfrac12\overline F_{\chi_2^2}(t)+
\tfrac12\overline F_{\chi_3^2}(t)$ because
$\overline F_{\chi_1^2}(t)\le
\overline F_{\chi_3^2}(t)$. For equality, the pathwise
nested-cone distance bound and the same wedge dominate its
Gaussian statistic. The preceding theorem transfers both
bounds uniformly to the likelihood ratios.
\end{proof}

\paragraph{Scope and numerical approximation.}
The fixed-stratum angular laws may be used where their
specific tangent-model assumptions hold; the uniform
conservative cutoffs above require no covariance-rank
selection. The sharper equality envelope uses the
\emph{exact-chart} information, not a sample Hessian evaluated
at an incorrectly imposed vertex. All three calibrations
concern ideal marginal likelihood suprema and exact nonlinear
null constraints. A numerical implementation inherits them
only if the full and null maximized log likelihoods, information
when used, and numerical critical values have the appropriate
uniform negligible errors. Those software requirements have
not been established by the theoretical proof.
